\paragraph{Memory helps, and more so the more correlated the noise.}
Once $\rho\gtrsim0.5$, every decoder in the parameter sweeps (E2--E4) beats
nearest lattice, by up to two orders of magnitude at $\rho\approx0.9$. Badly
mis-specified filters (E5) are the exception. The gain follows the
predictability of the next displacement, $\sigma_q\sqrt{1-\rho^2}$
(Section~\ref{sec:why}). Much of it, though not most, comes from being able
to \emph{revise} earlier decisions. At the nominal point, predicting the next
displacement alone gives a factor of about 13 for Bayes (about 9 for the
LSTM), and revision adds a further factor of about 7 (about 2.7 for the LSTM).

\paragraph{When the model is right, Bayes wins.}
With the correct noise model, the grid filter is the optimal syndrome-only
decoder for this action set, and it was best at every test point. The best
LSTMs come within 2--17\% of it where logical errors are common (matched
models at $\rho=0.5$: 2\%; at $\sigma_q=0.26L$: 17\%). They fall 3--4 times
behind where errors are rare (the nominal point): the decisive rounds are
rare in the training data, and the network has not learned them.

\paragraph{When the model is moderately wrong, Bayes usually still wins.}
A filter whose assumed parameters are moderately off loses little: every
model with $0.8\le\rho'\le0.95$ and $0.75\le\sigma_q'/\sigma_q\le1.5$ stays
within $1.6\times$ of the matched error. Overestimating $\sigma_q$ is cheap.
Underestimating the innovation noise, lowering $\rho'$ substantially, or
ignoring the correlation is expensive, and can be worse than no memory at all
(E5).

In the parameter sweeps, the fixed nominal filter stayed close to matched
Bayes over a wide range of $\sigma_q$ and for $\rho\ge0.8$. Averaged over
LSTM seeds, its error was at or below that of the nominally trained LSTM at
every sweep point; at $\rho=0$ the two were within the LSTM seed spread.

A Bayesian average over a coarse grid of models (Bayes-broad) matched or beat
the broadly trained LSTM at most test points, including most of level~B,
where both are out of distribution. The exceptions are weak-to-moderate
correlation inside the training box: the LSTM was better at level~A with
$\rho=0.3$ for all three seeds, and slightly better at $\rho=0.5$.

\paragraph{When the model is structurally wrong, ML helps.}
Under regime-switching correlation (E6), every fixed-model filter is
mis-specified, and the LSTM trained on varied noise beat all of them. Only a
filter that models the switching did better. This is the clearest support in
our data for the proposal's hypothesis that machine learning helps when the
Bayesian model is inaccurate. The advantage came from training on varied
correlations rather than on the switching process itself.

\paragraph{Generalisation needs varied training.}
\begin{itemize}
  \item A network trained at one point degrades away from it. At $\rho=0$ it
    is worse than NL on average over seeds, because it applies memory that is
    not there.
  \item Training on a broad box removes most of that failure, at the cost of
    some in-distribution performance and a larger spread between seeds.
  \item The same holds for Bayes: a single assumed model is fragile in the
    same places, and averaging over models removes most of the problem. At
    $\rho=0$, Bayes-broad is still slightly worse than NL.
\end{itemize}

\paragraph{Cost.}
On a GPU, both memory decoders process about $3\times10^5$ trial-rounds per
second, with a latency of a few milliseconds per round. On a CPU, the
single-trial latency is similar (about 0.9~ms per round for both), but for
batched decoding the LSTM's throughput is about 27 times the grid filter's.
The filter's cost grows with the number of hidden variables, which favours
learned decoders for richer noise models. The LSTM, in turn, needs training
data and about 6--11 minutes of GPU time per model.

\paragraph{Answer to the research question.}
In this setting, machine learning did not outperform a correctly specified
Bayesian filter, and it rarely outperformed a mildly mis-specified one. The
exception is a model-averaged filter at weak correlation. A broadly trained
LSTM is nevertheless a robust, cheap and assumption-light decoder:
\begin{itemize}
  \item it beats nearest-lattice decoding by large factors under correlated
    noise;
  \item it beats mis-specified filters when the true noise is structurally
    different from what the filter assumes.
\end{itemize}
